\section{Marco teórico y trabajos relacionados}

\subsection{La Gestión de Talento Humano (GTH) en el entorno digital}
La GTH comprende los procesos de atracción, incorporación, desarrollo, compensación y retención del talento. En consecuencia, las soluciones digitales impulsan el seguimiento de los indicadores y contribuyen a una reducción en tiempos de procesos repetitivos. Algunos estudios revelan cómo el reclutamiento de personal, la formación y las mejoras en el desempeño se incrementan teniendo como punto de partida la interconexión entre plataformas y datos \cite{gestion_talento_digital}. En el caso de las pymes, el impacto es mayor debido a la no existencia de sistemas legados por defecto y la llegada de la posibilidad de implantar procesos estándar.

\subsection{Las modernas arquitecturas web}
Las aplicaciones web de la empresa han dado un giro hacia lo que son las arquitecturas de servicios existentes, con capas claramente definidas (interfaz de presentación, lógica y datos). En esta línea, el Spring Framework y el ecosistema que lo envuelve (Spring Boot, Spring Security, Spring Data...) son los que dominan el panorama de productividad y mantenibilidad en el mundo de Java, al igual que el soporte nativo que proporciona la seguridad, la gestión de datos o las pruebas de software \cite{estudio_aplicaciones_java}\cite{analisis_comparativo_java}. React, en contraposición, permitiendo estructurar la interfaz de usuario mediante componentes y un ciclo de vida declarativo, favorece la reusabilidad y testabilidad de la misma \cite{analisis_frameworks_frontend}.

\subsection{Diseño responsivo y experiencia de usuario}
Bootstrap proporciona un sistema de rejillas, utilidades y componentes que simplifican y agilizan la construcción de interfaces con un grado de consistencia en distintas dimensiones de pantalla; al mismo tiempo que huimos de la aparición de deuda técnica e incrementamos la accesibilidad \cite{optimizacion_bootstrap}. La evidencia sugiere que los sistemas adoptados con la interfaz de usuario con una misma estructura, con la existencia de un mismo feedback claro, traen consigo el incremento del uso de la misma.

\subsection{Modelado de datos relacional}
Las claves primarias y foráneas, la normalización y el control de integridad permiten garantizar la calidad de los datos y evitar las anomalías de actualización. La literatura indica que un buen diseño conceptual y lógico reduce los costes de mantenimiento y favorece la analítica posterior \cite{analisis_google_charts}.

\subsection{Visualización y analítica operacional}
Las bibliotecas Google Charts permiten construir paneles de control que miden la asistencia, los permisos y la rotación de personas, además la integración de estos con fuentes relacionales es un patrón común en los mismos tableros operativos \cite{analisis_google_charts}.

\subsection{Microservicios y estandarización}
La estandarización de contratos de servicio, la observabilidad y el versionado son esenciales durante la migración a microservicios. Informes de casos de estudio evidencian mejoras en latencias para servicios de lectura de gran carga mediante la especialización de los mismos \cite{implementacion_microservicios}. En pymes, una traslación progresiva es más recomendable puesto que limita los riesgos y los costes.

\subsection{Competencias técnicas y ecosistema de lenguajes}
La elección del lenguaje de programación viene respaldada por el mercado (Java y JavaScript son soportados por la demanda, se reproduce la disponibilidad de proyectos y librerías y también la comunidad, podría justificar su uso en proyectos académicos o empresariales \cite{analisis_lenguajes}). React Native y Flutter son alternativas para construcciones móviles, sin embargo, para intranets y flujos administrativos, una SPA web permitirá tener la menor fricción de despliegue \cite{react_native_distancias}.